\chapter{Toward Quaternionic Separation}
\label{ch:quaternionic}

Chapter~\ref{ch:minimal-witnesses}'s second open question,
$m_\beta(\C,\Hset;d)$, was left entirely untouched pending the
composite-system chapters. Those chapters (Ch.~\ref{ch:local-tomography}--
\ref{ch:entanglement}) found quaternionic composition genuinely
obstructed — the naive product of local quaternionic Hermitian bases is
already overcomplete (Ch.~\ref{ch:tensor-product}'s remark). This chapter
asks a narrower question first: not composite systems, but whether the
single-system graph-realizability machinery of
Chapters~\ref{ch:context-graphs}--\ref{ch:minimal-witnesses} — which never
required composing separate systems, only composing frame transitions of
one system along a graph — extends to $\Hset$ at all.

\section{What Generalizes, and Why}

The graph-realizability machinery (Proposition~\ref{prop:tree-no-constraint},
Theorem~\ref{thm:graph-realizability}) was built entirely from matrix
\emph{multiplication} of frame transitions: $F_v := F_uL_{uv}$ propagated
along a tree, and cycle closure $L_{A_0A_1}\cdots L_{A_{k-1}A_0}=I$ via
telescoping. Neither proof used commutativity anywhere — only
associativity, and the fact that $\pi(L)_{ij}$ (the modulus-squared of a
matrix entry) is well-defined and nonnegative.

\begin{proposition}[The machinery is field-agnostic]
\label{prop:h-generalizes}
\index{The machinery is field-agnostic (proposition)}
Proposition~\ref{prop:tree-no-constraint} and
Theorem~\ref{thm:graph-realizability} hold verbatim with $\R$ replaced by
$\Hset$ throughout, where $\pi(L)_{ij} := |L_{ij}|_\Hset^2 = L_{ij}
\overline{L_{ij}}$ (the quaternionic norm, a nonnegative real number for
every quaternion $L_{ij} = a+bi+cj+dk$), and "unitary" means
$L^\dagger L = I$ with $\dagger$ the quaternionic conjugate transpose.
\end{proposition}

\begin{proof}
Quaternionic matrix multiplication is associative (inherited from
associativity of quaternion multiplication itself) though not commutative.
Every step of Proposition~\ref{prop:tree-no-constraint}'s proof (unique
tree-path propagation) and Theorem~\ref{thm:graph-realizability}'s proof
(construct from a tree-lift choice, check non-tree edges) used only
associativity and the definition of $\pi$; neither invoked commutativity of
matrix entries or of the underlying scalars.
\end{proof}

This was checked directly, not merely asserted: quaternionic associativity,
a genuine quaternionic unitary matrix (via quaternionic Gram--Schmidt)
satisfying $U^\dagger U = I$, exact cycle closure $L_{ZX}L_{XY}L_{YZ}=I$ for
random quaternionic frames, and exact reproduction of tree-edge data under
propagation, all confirmed numerically before this proposition was written.

\section{The \texorpdfstring{$\C$--$\Hset$}{C-H} Minimal Witness Question}

Definitions~\ref{def:realizability-space}--\ref{def:minimal-witness} are
already stated for general $\mathbb F$, so $m_\beta(\C,\Hset;d)$ is
well-posed immediately given
Proposition~\ref{prop:h-generalizes}. Two facts constrain where a witness
could live.

\begin{proposition}[$\C \subset \Hset$ as an inclusion of admissible sets]
\label{prop:c-subset-h}
\index{C subset H as an inclusion of admissible sets (proposition)}
$\mathcal R_\C(G,d) \subseteq \mathcal R_\Hset(G,d)$ for every $G,d$.
\end{proposition}

\begin{proof}
$\C$ embeds as a subalgebra of $\Hset$ ($a+bi \hookrightarrow a+bi+0j+0k$),
so every complex unitary matrix is a quaternionic unitary matrix with
entries restricted to this subalgebra; any real frame assignment achieving
a labeling over $\C$ achieves it over $\Hset$ using the identical matrices.
\end{proof}

\begin{proposition}[No single-edge witness at $d=2$]
\label{prop:no-d2-witness}
\index{No single-edge witness at d=2 (proposition)}
$\mathcal R_\R(G,2) = \mathcal R_\C(G,2) = \mathcal R_\Hset(G,2)$ for every
tree $G$; in particular $m_\beta(\C,\Hset;2) \ge 1$, if it exists at all.
\end{proposition}

\begin{proof}
Theorem~\ref{thm:n2degeneracy} shows every $B\in\mathrm{DS}(2)$ is already
orthostochastic: $\mathrm{OS}(2)=\mathrm{DS}(2)$. Since $\R\subset\C\subset
\Hset$, the same real orthogonal matrix witnessing $\R$-admissibility
witnesses $\C$- and $\Hset$-admissibility too, so all three edgewise
admissible sets coincide at $d=2$, and by
Proposition~\ref{prop:tree-no-constraint} (both real and quaternionic
versions) no tree can separate any pair of the three fields at $d=2$.
\end{proof}

\section{What This Chapter Does Not Establish}

Unlike Chapter~\ref{ch:higher-d-witnesses}'s resolution of the real/complex
question at $d\ge3$, no analogous collapse is established here. That
resolution rested on an explicit, verified single matrix
(Theorem~\ref{thm:n3incompleteness}'s DFT-derived $\Gamma \in
\mathrm{US}(3)\setminus\mathrm{OS}(3)$) exhibiting the separation directly.
The analogous question for $\C$ versus $\Hset$ — does there exist, at some
small $d$, a quaternionic-unistochastic matrix admitting no complex
realization — has no such explicit witness in hand. Constructing one
requires either an explicit small-dimension quaternionic unitary matrix
whose modulus-squared provably escapes every complex realization (the
direct analogue of the DFT construction, but genuinely harder: the
relevant non-existence fact would concern complex, not real, matrix
constraints, and no argument as elementary as the Hadamard-matrix
non-existence check of Theorem~\ref{thm:n3incompleteness} is currently in
hand), or a correct quaternionic generalization of the mutually-unbiased-
basis counting bound that underlies Chapter~\ref{ch:real-mub-obstruction}'s
triangle argument — and the standard derivation of that bound leans on
structural facts about the complex traceless-Hermitian-operator space that
do not obviously carry over: a quaternionic orthonormal basis carries
$\mathrm{Sp}(1)$ (three real parameters) of local phase freedom per vector
rather than the complex case's $U(1)$ (one real parameter), which changes
the dimension count in a way this book has not worked out correctly.
Rather than risk asserting an unverified bound, this is left as the
precise, well-posed target for future work.

\begin{remark}[The gap, isolated]
\label{rem:gap-isolated}
Propositions~\ref{prop:c-subset-h} and~\ref{prop:no-d2-witness} narrow
$m_\beta(\C,\Hset;d)$, if it exists at all, to $d\ge3$. A separation at
some such $d$ would follow from either:
\begin{enumerate}
\item[(i)] an explicit bistochastic matrix on some small $d\ge3$ realizable
quaternionically but not complexly — the direct analogue of
Theorem~\ref{thm:n3incompleteness}'s DFT construction, requiring a
non-existence fact about complex matrices in place of that theorem's
non-existence of a real Hadamard matrix; or
\item[(ii)] a structural obstruction — most plausibly a corrected
quaternionic mutually-unbiased-basis counting bound, redone from first
principles to account for the $\mathrm{Sp}(1)$ (rather than $U(1)$) local
phase freedom per basis vector — proving such a matrix must exist without
exhibiting one directly, in the style of
Chapter~\ref{ch:real-mub-obstruction}'s cycle argument.
\end{enumerate}
Neither is completed here. What this chapter contributes is the
reduction: the search is not open at every $d$, only at $d\ge3$, and the
machinery needed to verify a candidate witness, once found, is already in
place (Proposition~\ref{prop:h-generalizes}).
\end{remark}

\begin{ledgerentry}[End of Chapter~\ref{ch:quaternionic}]
\textbf{Established:} the graph-realizability machinery of
Chapters~\ref{ch:context-graphs}--\ref{ch:minimal-witnesses} generalizes to
$\Hset$ verbatim, verified computationally
(Proposition~\ref{prop:h-generalizes}); $\mathcal R_\C \subseteq \mathcal
R_\Hset$ always (Proposition~\ref{prop:c-subset-h}); no single-edge witness
exists at $d=2$ for any pairing of $\R,\C,\Hset$
(Proposition~\ref{prop:no-d2-witness}).\\
\textbf{Not established:} $m_\beta(\C,\Hset;d)$ for any $d$ — neither an
explicit witness nor a correct counting bound. This chapter narrows the
search (rules out $d=2$ single edges, confirms the machinery needed to
search $d\ge3$ is sound) without completing it. The quaternionic
MUB-counting question, and the composite-system obstruction flagged in
Chapter~\ref{ch:tensor-product}, remain the two live approaches.
\end{ledgerentry}
